\documentclass[11pt,reqno]{amsart}
\usepackage[T1]{fontenc}
\usepackage{lmodern}
\usepackage{microtype}
\usepackage{amsmath,amssymb,mathtools}
\usepackage[colorlinks=true,linkcolor=blue,citecolor=blue,urlcolor=blue]{hyperref}

\newtheorem{theorem}{Theorem}[section]
\newtheorem{lemma}[theorem]{Lemma}
\newtheorem{proposition}[theorem]{Proposition}
\newtheorem{corollary}[theorem]{Corollary}
\newtheorem{conjecture}[theorem]{Conjecture}
\theoremstyle{definition}
\newtheorem{definition}[theorem]{Definition}
\newtheorem{remark}[theorem]{Remark}
\newtheorem{example}[theorem]{Example}

\newcommand{\E}{\mathbb E}
\newcommand{\PP}{\mathbb P}
\newcommand{\supp}{\operatorname{supp}}
\newcommand{\energy}{\operatorname{En}}
\newcommand{\Ho}{\mathcal H}
\newcommand{\cE}{\mathcal E}
\newcommand{\cJ}{\mathcal J}
\newcommand{\cM}{\mathcal M}
\newcommand{\cC}{\mathcal C}
\newcommand{\cD}{\mathcal D}
\newcommand{\cP}{\mathcal P}
\newcommand{\Lc}{\mathcal L}
\raggedbottom

% TODO(submission): replace "Anonymous" by the author metadata (kept anonymous for internal circulation).
\title[An energy bound for avoidance indicators]{An energy bound for avoidance indicators,\\
and Fourier tails of DNFs over product spaces}
\author{Anonymous}
\date{Draft v1, 5 October 2026}

\begin{document}

\begin{abstract}
Let $F$ be the indicator that none of a finite family of cylinder events
occurs, on an arbitrary finite product probability space, and let
$F=\sum_UF^{=U}$ be its Efron--Stein decomposition.  We prove that
\[
 \sum_U\Big(\prod_{v\in U}\lambda_v\Big)\|F^{=U}\|_2^2\ \le\ 1
\]
for all weights $\lambda_v\ge1$ such that $\prod_{v\in\supp E}\lambda_v\le2$ for every
event $E$.  There is no dependence on the number of events, the alphabet
sizes or the measure.  In particular every width-$k$ DNF
$g:\{\pm1\}^n\to\{\pm1\}$ satisfies $\sum_{|S|>t}\hat g(S)^2\le4\cdot2^{-(t+1)/k}$, under
every product measure; so $g$ is $\varepsilon$-concentrated up to degree
$k\log_2(4/\varepsilon)$.  The proof does not use a switching lemma.  It bounds the
Efron--Stein energies by signed cover counts of the events that hold at a
random point, writes the resulting hypergraph functional as the second moment
of an alternating union sum under a random deletion, and bounds that sum by
deletion--contraction.  Over general product spaces the base $2^{-1/k}$ is
sharp, and the threshold $2$ is necessary.  We also give a version with
weights that grow along a filtration of each coordinate, and we explain the
application that motivated it: the junta term in a lower bound for the
least ``multiplier modulus'' of the Erd\H{o}s--Straus equation.
\end{abstract}

\maketitle
% Work around a pdfTeX/microtype font-expansion warning (as in the other notes).
\setbox0\hbox{\footnotesize\textsuperscript{1}}

\section{Introduction}\label{sec:intro}

Let $\Omega=\prod_{v\in V_0}\Omega_v$ be a finite product of finite sets with a
product probability measure $\pi=\bigotimes_v\pi_v$.  A \emph{cylinder event}
fixes the values of the coordinates in a set $\supp E\subseteq V_0$, its
\emph{support}.  For a finite family $\cE$ of cylinder events let
\[
 F=F_\cE=\prod_{E\in\cE}(1-\mathbf 1_E)
\]
be its \emph{avoidance indicator}, and let $F=\sum_{U\subseteq V_0}F^{=U}$ be the
Efron--Stein (Hoeffding) decomposition of $F$ in $L^2(\pi)$
\cite{ES,Hoeffding}, \cite[\S8.3]{OD}.  For the uniform measure on
$\{\pm1\}^n$ this is the Fourier expansion, $F^{=U}=\hat F(U)\chi_U$.  We
study the weighted energy
\[
 G_F(\lambda)=\sum_{U\subseteq V_0}\lambda^U\|F^{=U}\|^2,\qquad \lambda^U=\prod_{v\in U}\lambda_v,\quad\lambda_v\ge1 .
\]
For constant weights $\lambda_v=\rho^2$ it is $\|T_\rho F\|^2$, where $T_\rho$
multiplies the component $F^{=U}$ by $\rho^{|U|}$; so we are asking for the
largest noise operator \emph{with $\rho>1$} that is still a contraction from
avoidance indicators to $L^2$.

\begin{theorem}[Energy bound]\label{thm:main}
Let $\lambda_v\ge1$ $(v\in V_0)$, and suppose that $\lambda^{\supp E}\le2$ for every
$E\in\cE$.  Then
\[
 G_F(\lambda)=\sum_U\lambda^U\|F^{=U}\|^2\ \le\ 1 .
\]
\end{theorem}

The hypothesis is on single events only.  There is no condition on the
number of events, on their masses or overlaps, on the alphabet sizes
$|\Omega_v|$ or on the measure.  The constant $2$ cannot be raised
(Proposition~\ref{prop:single}).  Since
$\sum_{U:\lambda^U>e^\tau}\|F^{=U}\|^2\le e^{-\tau}G_F(\lambda)$, the theorem is a family of
tail bounds, one for each choice of weights.  Constant weights give the
following.

\begin{corollary}[Energy tail]\label{cor:tail}
If every event of $\cE$ has $|\supp E|\le k$ $(k\ge1)$, then for
$f\in\{F,1-F\}$ and every integer $t\ge0$,
\[
 \energy(f;t):=\sum_{|U|>t}\|f^{=U}\|^2\ \le\ 2^{-(t+1)/k}\,\PP(F=0)\big(2-\PP(F=0)\big)\ \le\ 2^{-(t+1)/k}.
\]
\end{corollary}

A width-$k$ DNF is $g=2F-1$ for the family of its terms (with the convention
$-1=\mathrm{True}$), and the corollary becomes a statement about its Fourier
tail.

\begin{corollary}[Fourier tails of DNFs]\label{cor:dnf}
Let $g:\{\pm1\}^n\to\{\pm1\}$ be computed by a DNF of width $k$, and let $\pi$ be
any product measure on $\{\pm1\}^n$ (Efron--Stein components in $L^2(\pi)$;
Fourier coefficients if $\pi$ is uniform).  Then for every integer $t\ge0$
\[
 W^{>t}[g]:=\sum_{|U|>t}\|g^{=U}\|^2\ \le\ 4\,p(2-p)\cdot2^{-(t+1)/k}\ \le\ 4\cdot2^{-(t+1)/k},\qquad p=\PP[g=-1].
\]
So $g$ is $\varepsilon$-concentrated up to degree $k\log_2(4/\varepsilon)$.  Moreover
$\sum_U2^{|U|/k}\|g^{=U}\|^2\le1+4p$, and the total influence is at most
$(4k/\ln2)\,p$.  The same holds for $q$-ary DNFs, with literals $x_v\in S_v$,
under any product measure on $\prod_v[q_v]$.
\end{corollary}

\emph{Comparison.}  The classical bound is due to Linial, Mansour and Nisan
\cite{LMN}, via H{\aa}stad's switching lemma \cite{Hastad}: width-$w$ DNFs are
$\varepsilon$-concentrated up to degree $O(w\log(1/\varepsilon))$.  In the form of
O'Donnell's book \cite[\S4.4]{OD} it reads $W^{\ge t}[g]\le2\cdot2^{-t/(20w)}$
(we quote the constant $20$ from memory and could not re-check it).  It has
been used in this form, with an unspecified constant $C$, in later work on
DNFs, for instance by Mansour \cite{Mansour} and by Lecomte and Tan
\cite[Fact~6]{LT}.  Corollary~\ref{cor:dnf} has the constant $1$ in the
exponent and holds for arbitrary product measures and alphabets.  Over
general product spaces the base $2^{-1/k}$ is sharp: no bound
$C\rho^{-t/k}$ with $\rho>2$ holds, and at most a factor of order $\sqrt{t/k}$
can be gained (Proposition~\ref{prop:sharp}).  On the uniform cube we only
know that the true rate lies between $3^{-t/k}$ and $2^{-t/k}$
(Example~\ref{ex:parity}).  The bound on the total influence is weaker by a
constant factor than the known bound $2w$ for width-$w$ DNFs (we recall this
from memory; it is usually attributed to Boppana \cite{Boppana}), and we
state it only because it falls out.

The nearest prior art we know is the paper of Lecomte and Tan \cite{LT}.
They bound $|\hat g(S)|$ by $2^{|S|}$ times the probability that $S$ is covered
by the terms satisfied at a random point.  This is the device of our
Lemma~\ref{lem:cover}, but there are three differences: our cover counts are
\emph{signed}, and vanish as soon as a satisfied term misses $S$; we work in
$L^2$ with all levels and weights at once; and our degree concentration does
not go through the switching lemma.  The two combinatorial steps of the proof,
a polarization identity (Lemma~\ref{lem:polar}) and a deletion--contraction
bound (Lemma~\ref{lem:theta}), have no counterpart there.

\emph{What we checked.}  We have no internet access while writing this
draft, and our literature search is partial.  Earlier, \cite{LT} and Tal
\cite{Tal} were checked online; in \cite{Tal} we found no explicit-constant
tail bound for width-$w$ DNFs.  Everything else above is quoted from memory and is
marked as such.  We did not search systematically through the switching-lemma
literature (sharp variants of H{\aa}stad's lemma, Rossman's entropy proof),
the literature on Fourier growth, or work on noise operators with $\rho>1$.
We believe Theorem~\ref{thm:main} and the constant $1$ in
Corollary~\ref{cor:dnf} are new, but we make no priority claim.

\emph{Proof outline.}  Write $G_F=\sum_V\mu^V\|L_VF\|^2$ with $\mu_v=\lambda_v-1$ and
$L_V=\prod_{v\in V}(I-\E_v)$.  Section~\ref{sec:cover} bounds $\|L_VF\|^2$ by the
mean square of a signed count $N_{\Ho(x)}(V)$ of covers of $V$ by the events
that hold at $x$, and so reduces Theorem~\ref{thm:main} to the hypergraph
inequality $Q_\mu(\Ho)=\sum_V\mu^VN_\Ho(V)^2\le1$.  Section~\ref{sec:hyper} proves
it: $Q_\mu(\Ho)$ is the second moment of the alternating sum
$\Theta_\lambda(\Ho_P)=\sum_{\cJ\subseteq\Ho_P}(-1)^{|\cJ|}\lambda^{\bigcup\cJ}$ over the edges that
miss a random set $P$, and deletion--contraction gives $|\Theta|\le1$.  Nothing
in the proof is specific to two-point spaces or to the uniform measure.

\emph{Organisation.}  Section~\ref{sec:prelim} sets up notation.
Sections~\ref{sec:cover} and~\ref{sec:hyper} prove Theorem~\ref{thm:main}.
Section~\ref{sec:dnf} proves the corollaries, Section~\ref{sec:sharp} discusses
sharpness, failed strengthenings and open problems, and
Section~\ref{sec:filtration} gives the filtration variant.
Section~\ref{sec:app} sketches the application to the Erd\H{o}s--Straus
equation \cite{Subexp}.  The exact-arithmetic script
\texttt{scripts/energy\_note\_check.py} checks the main inequalities on small
cases (Section~\ref{sec:check}).

\section{Setting and elementary facts}\label{sec:prelim}

\emph{Product spaces.}  Throughout, $\Omega=\prod_{v\in V_0}\Omega_v$ is a finite product
of finite sets, $V_0$ is finite, and $\pi=\bigotimes_v\pi_v$ is a product
probability measure.  We write $\PP$, $\E$, $\langle\cdot,\cdot\rangle$ and
$\|\cdot\|=\|\cdot\|_{L^2(\pi)}$.  For $W\subseteq V_0$ we write $x_W$ for the restriction
of $x\in\Omega$ to $W$, $\Omega_W=\prod_{v\in W}\Omega_v$ and $\pi_W=\bigotimes_{v\in W}\pi_v$.  A
\emph{cylinder event} is a set $E=\{x:\ x_v=\sigma_v\ (v\in\supp E)\}$ with
$\supp E\subseteq V_0$ and $\sigma\in\Omega_{\supp E}$.  A family $\cE$ of cylinder events is
finite, and repetitions are allowed.  Its avoidance indicator is
$F=\prod_{E\in\cE}(1-\mathbf 1_E)$, and $h=1-F$ is the indicator that some event occurs.
For weights $\lambda_v\ge1$ we put $\mu_v=\lambda_v-1\ge0$, $\lambda^U=\prod_{v\in U}\lambda_v$,
$\mu^U=\prod_{v\in U}\mu_v$, and $w_E=\lambda^{\supp E}$ (the \emph{weight} of $E$).

\emph{Efron--Stein decomposition} (classical; \cite{ES,Hoeffding},
\cite[\S8.3]{OD}).  For $\varphi:\Omega\to\mathbb R$ and $W\subseteq V_0$ let
$\E[\varphi\mid x_W]$ be the conditional expectation given the coordinates in
$W$, and
\[
 \varphi^{=U}=\sum_{W\subseteq U}(-1)^{|U\setminus W|}\,\E[\varphi\mid x_W].
\]
Then $\varphi=\sum_U\varphi^{=U}$, the components are mutually orthogonal,
$\varphi^{=U}$ depends only on $x_U$, and $\E_u\varphi^{=U}=0$ for $u\in U$, where
$\E_u$ averages over the coordinate $u$ (with respect to $\pi_u$).  Put
$L_u=I-\E_u$ and $L_V=\prod_{u\in V}L_u$.  The $L_u$ are commuting orthogonal
projections, and $L_V\varphi=\sum_{U\supseteq V}\varphi^{=U}$.  If $\cD$ is a down-closed
family of subsets of $V_0$, then $\varphi^{\cD}=\sum_{U\in\cD}\varphi^{=U}$ is a linear
combination of the functions $\E[\varphi\mid x_W]$, $W\in\cD$ (by M\"obius
inversion), and
\begin{equation}\label{eq:energy}
 \E(\varphi-\varphi^{\cD})^2=\sum_{U\notin\cD}\|\varphi^{=U}\|^2 .
\end{equation}
For $\cD=\{U:\ |U|\le t\}$ this is the \emph{energy above level $t$},
$\energy(\varphi;t)$.  It is the least value of $\E(\varphi-g)^2$ over sums $g$
of functions that each depend on at most $t$ coordinates.  For $\pm1$-valued
$g$ we also write $W^{>t}[g]=\energy(g;t)$, as in \cite{OD}.

\begin{definition}[Energy generating function]
For weights $\lambda_v\ge1$,
\[
 G_\varphi(\lambda)=\sum_{U\subseteq V_0}\lambda^U\|\varphi^{=U}\|^2 .
\]
\end{definition}

\begin{lemma}[Elementary identities]\label{lem:ident}
For every $\varphi$ and all weights $\lambda_v\ge1$:
\begin{enumerate}
\item[(a)] $G_\varphi(\lambda)=\sum_{V\subseteq V_0}\mu^V\|L_V\varphi\|^2$;
\item[(b)] $\sum_{U:\ \lambda^U>e^\tau}\|\varphi^{=U}\|^2\le e^{-\tau}G_\varphi(\lambda)$ for every real $\tau$;
\item[(c)] $G_{1-\varphi}=1-2\E\varphi+G_\varphi$; in particular $G_F\le1$ if and only if
$G_h\le2\E h$;
\item[(d)] if $\varphi(x)=\prod_{v}\varphi_v(x_v)$ is a product, then
$G_\varphi(\lambda)=\prod_vG_{\varphi_v}(\lambda_v)$, where $G_{\varphi_v}(\lambda_v)=(\E\varphi_v)^2+\lambda_v\operatorname{Var}\varphi_v$.
\end{enumerate}
\end{lemma}

\begin{proof}
(a) Expand $\lambda^U=\prod_{v\in U}(1+\mu_v)=\sum_{V\subseteq U}\mu^V$ and use
$\|L_V\varphi\|^2=\sum_{U\supseteq V}\|\varphi^{=U}\|^2$.  (b) is Markov's inequality.
(c) $(1-\varphi)^{=U}=-\varphi^{=U}$ for $U\ne\emptyset$ and $(1-\varphi)^{=\emptyset}=1-\E\varphi$.
(d) For a product, $\varphi^{=U}=\prod_{v\in U}(\varphi_v-\E\varphi_v)\prod_{v\notin U}\E\varphi_v$.
\end{proof}

\begin{proposition}[A single event; the threshold $2$]\label{prop:single}
Let $E$ fix $x_v=\sigma_v$ on $\supp E$, put $p_v=\pi_v(\sigma_v)$ and
$\pi_E=\PP(E)=\prod_{v\in\supp E}p_v$.  Then
\[
 G_{1-\mathbf 1_E}(\lambda)=1-2\pi_E+\pi_E\prod_{v\in\supp E}\big(\lambda_v-\mu_vp_v\big).
\]
Hence, for fixed weights with $w_E>2$, $G_{1-\mathbf 1_E}(\lambda)>1$ as soon as all
$p_v>0$ are small enough.  So the constant $2$ in Theorem~\ref{thm:main}
cannot be replaced by any larger number, even for a single event.
\end{proposition}

\begin{proof}
By Lemma~\ref{lem:ident}(d), with $\varphi_v=\mathbf 1[x_v=\sigma_v]$,
$G_{\mathbf 1_E}=\prod_v(p_v^2+\lambda_vp_v(1-p_v))=\pi_E\prod_v(\lambda_v-\mu_vp_v)$.  Then use
Lemma~\ref{lem:ident}(c).  The product tends to $w_E$ as $p\to0$.
\end{proof}

On the uniform cube $\{\pm1\}^n$ the single-event condition is the weaker
$\prod_v(1+\lambda_v)/2\le2$.  The threshold $w_E\le2$ is forced by large
alphabets or strongly biased measures.

%%SECTIONS%%

\begin{thebibliography}{99}
\bibitem{Boppana} R.~B.~Boppana, The average sensitivity of bounded-depth
circuits, Inform.\ Process.\ Lett.\ 63 (1997), 257--261.
% TODO(verify): [Boppana], [Mansour], [Tal] details written from memory.
\bibitem{ES} B.~Efron and C.~Stein, The jackknife estimate of variance, Ann.\
Statist.\ 9 (1981), 586--596.
\bibitem{Hastad} J.~H{\aa}stad, Almost optimal lower bounds for small depth
circuits, in: Proc.\ 18th ACM STOC (1986), 6--20.
\bibitem{Hoeffding} W.~Hoeffding, A class of statistics with asymptotically
normal distribution, Ann.\ Math.\ Statist.\ 19 (1948), 293--325.
\bibitem{LT} V.~Lecomte and L.-Y.~Tan, Sharper bounds on the Fourier
concentration of DNFs, in: Proc.\ 62nd IEEE FOCS (2021), arXiv:2109.04525.
\bibitem{LMN} N.~Linial, Y.~Mansour and N.~Nisan, Constant depth circuits,
Fourier transform, and learnability, J.\ ACM 40 (1993), 607--620.
\bibitem{Mansour} Y.~Mansour, An $O(n^{\log\log n})$ learning algorithm for DNF
under the uniform distribution, J.\ Comput.\ System Sci.\ 50 (1995), 543--550.
\bibitem{OD} R.~O'Donnell, Analysis of Boolean Functions, Cambridge University
Press, 2014.
\bibitem{Subexp} Anonymous, Large Erd\H{o}s--Straus witnesses II: the multiplier
modulus is infinitely often $\exp((\log p)^{1/5-o(1)})$, preprint, 2026.
\bibitem{Tal} A.~Tal, Tight bounds on the Fourier spectrum of $\mathrm{AC}^0$, in:
Proc.\ 32nd CCC (2017), LIPIcs 79, Art.~15; ECCC TR14-174.
\end{thebibliography}

\end{document}
